Fix a sufficiently large triangular-array index at which \(\mathcal M_{n,d,L}\neq\varnothing\), and select the kernel and spectral data of any one member.  The proof first constructs from those data a graph marginal which is admissible for the same parameter tuple and whose degrees have a distribution independent of the selected kernel.

\textit{Step 1: a null-row Erd\H{o}s--R\'enyi embedding.}
Let \(T_W\) be the integral operator with kernel \(W\), and set
\[
 b_k=\int_0^1\psi_k(u)\,du,
 \qquad 1\leq k\leq r.
\]
The functions \(\psi_k\) are integrable by Cauchy--Schwarz and their \(L^2\)-normalization in \({\rm ass:eigenfunction\text{-}orthonormality}\).  The finite expansion in \({\rm ass:finite\text{-}rank}\), Fubini's theorem for a finite sum, and \({\rm ass:degree\text{-}regularity}\) give
\[
 \sum_{k=1}^r\lambda_k b_k^2
 =\int_0^1\!\int_0^1W(u,v)\,du\,dv
 \geq\kappa^{-1}>0.
\tag{1}
\]
Consequently at least one eigenvalue is positive; fix \(k_+\) with \(\lambda_{k_+}>0\).

Choose one point \(u_0\in[0,1]\).  Replacing representatives on a singleton, define
\[
 \psi_k^*(u)=\psi_k(u)\quad(u\neq u_0),
 \qquad
 \psi_{k_+}^*(u_0)=\lambda_{k_+}^{-1/2},
 \qquad
 \psi_k^*(u_0)=0\quad(k\neq k_+),
\tag{2}
\]
and define the everywhere-specified symmetric representative
\[
 W^*(u,v)=\sum_{k=1}^r\lambda_k\psi_k^*(u)\psi_k^*(v).
\tag{3}
\]
Equations (2)--(3) change the original spectral representative only on
\(\{u_0\}\times[0,1]\cup[0,1]\times\{u_0\}\), a set of two-dimensional Lebesgue measure zero.  Hence the \(L^2\) integral operator, its exact rank, its nonzero eigenvalues, orthonormality, the almost-everywhere kernel bound in \({\rm ass:kernel\text{-}bound}\), and the almost-everywhere row-mass inequality in \({\rm ass:degree\text{-}regularity}\) are unchanged.  In particular, \({\rm ass:spectral\text{-}strength}\) is unchanged.  At the selected point,
\[
 W^*(u_0,u_0)=\lambda_{k_+}\lambda_{k_+}^{-1}=1.
\tag{4}
\]

Now take \(U_i=u_0\) almost surely, independently across \(i\).  Each \(\psi_k^*(U_i)\) is constant, so its centered moments of every positive order are zero.  It therefore satisfies Bernstein's moment inequality with the prescribed parameter \(b_\psi>0\), verifying \({\rm ass:eigenfunction\text{-}bernstein}\).  Conditional on these latent types, draw the edges according to \({\rm ass:graphon\text{-}draw}\).  By (4), for \(i<j\),
\[
 A_{ij}\overset{\mathrm{ind}}{\sim}{\rm Bernoulli}(p_n),
 \qquad p_n=\rho_nW^*(u_0,u_0)=d/n.
\tag{5}
\]
The probability in (5) belongs to \((0,1)\) because the symbol declaration has \(0<d<n\); it is also consistent with \({\rm ass:edge\text{-}probability}\).  Thus (2)--(5), together with independent assignment from \({\rm ass:bernoulli\text{-}design}\), satisfy every graph, sampling, rank, support, and design property in \({\rm def:model\text{-}class}\).  Notice that this argument uses only nonemptiness to obtain one admissible integral operator.  It does not require either compatibility inequality of the block witness.

\textit{Step 2: a common two-point distance bound.}
Put \(s_0=B/4\) and let the i.i.d. variables \(\xi_i\), independent of \((U,A,Z)\), have density
\[
 p_{s_0}(t)=s_0^{-1}\cos^2\!\left(\frac{\pi t}{2s_0}\right)
 \mathbf1\{|t|\leq s_0\}.
\tag{6}
\]
The integral of (6) is one.  For a deterministic \(g\in C^3[0,1]\), define two response laws by
\[
 f_i^\pm(a,x)=\xi_i\pm g(x),
 \qquad a\in\{0,1\},\quad x\in[0,1].
\tag{7}
\]
Whenever
\[
 \|g\|_\infty\leq 3B/4,
 \qquad
 \max_{1\leq j\leq2}\|g^{(j)}\|_\infty\leq B,
 \qquad
 \|g^{(3)}\|_\infty\leq L,
\tag{8}
\]
the bound \(|\xi_i|\leq B/4\), (7), and (8) verify \({\rm ass:response\text{-}envelope}\) and \({\rm ass:third\text{-}derivative\text{-}bound}\).  The response functions are i.i.d. jointly with the constant latent types, and \({\rm ass:anonymous\text{-}interference}\) defines their potential outcomes.  Thus both signs in (7) are laws in the original class.  By \({\rm def:target}\), their target values are deterministic and satisfy
\[
 \theta_n^\pm=\pm g'(1/2),
 \qquad
 \Delta_g:=|\theta_n^+-\theta_n^-|=2|g'(1/2)|.
\tag{9}
\]

Let \(q_{s_0}=\sqrt{p_{s_0}}\), extended by zero outside \([-s_0,s_0]\).  It is absolutely continuous because it vanishes at the endpoints, and direct integration gives
\[
 \|q_{s_0}'\|_2^2
 =\frac{\pi^2}{4s_0^2}.
\tag{10}
\]
For translations \(\eta_1,\eta_2\), the fundamental theorem of calculus and Minkowski's integral inequality therefore imply
\[
 \begin{aligned}
 H^2\{p_{s_0}(\mathord\cdot-\eta_1),p_{s_0}(\mathord\cdot-\eta_2)\}
 &=\|q_{s_0}(\mathord\cdot-\eta_1)-q_{s_0}(\mathord\cdot-\eta_2)\|_2^2\\
 &\leq (\eta_1-\eta_2)^2\|q_{s_0}'\|_2^2
 =\frac{\pi^2(\eta_1-\eta_2)^2}{4s_0^2}.
\end{aligned}
\tag{11}
\]
Here \(H^2\) denotes the integral of the squared difference of square-root densities.  Define the observed exposure
\[
 X_i=\begin{cases}1/2+V_i,&N_i>0,\\0,&N_i=0.\end{cases}
\tag{12}
\]
Conditional on \((A,Z)\), the observed outcomes under (7) are independent translates of (6), with the two translations separated by \(2g(X_i)\).  The product-affinity identity, the inequality \(1-\prod_i(1-x_i)\leq\sum_i x_i\) for \(x_i\in[0,1]\), and (11) imply for the two complete observed-data laws \(Q_g^+,Q_g^-\) that
\[
 H^2(Q_g^+,Q_g^-)
 \leq \frac{\pi^2}{s_0^2}\mathbb E\sum_{i=1}^ng(X_i)^2.
\tag{13}
\]
The graph and assignment marginals are common to the two laws, which justifies averaging the conditional bound.  Cauchy--Schwarz applied to the two square-root densities gives
\[
 \operatorname{TV}(Q_g^+,Q_g^-)\leq H(Q_g^+,Q_g^-).
\tag{14}
\]

\textit{Step 3: degree and lattice estimates under the embedded graph.}
By (5), for every \(i\),
\[
 N_i\sim{\rm Binomial}(n-1,p_n),
 \qquad
 \mu_n:=\mathbb EN_i=(n-1)d/n.
\tag{15}
\]
For all sufficiently large \(n\), \(d/2\leq\mu_n\leq d\).  Conditional on \(N_i=N>0\), Bernoulli assignment gives \(M_i\sim{\rm Binomial}(N,1/2)\), and hence
\[
 \mathbb E_Z(V_i^2\mid A,N_i=N)=\frac1{4N}.
\tag{16}
\]
The binomial identity
\[
 \mathbb E\frac1{N_i+1}
 =\frac{1-(1-p_n)^n}{np_n}\leq\frac1d
\tag{17}
\]
and \(N^{-1}\leq2(N+1)^{-1}\) for \(N\geq1\) yield
\[
 \mathbb E\sum_{i=1}^nV_i^2\mathbf1\{N_i>0\}
 \leq\frac{n}{2d}.
\tag{18}
\]
Moreover,
\[
 n\Pr(N_i=0)=n(1-d/n)^{n-1}
 \leq ne^{-\mu_n}\leq ne^{-d/2}.
\tag{19}
\]

We also need a local binomial-mass estimate.  The largest atom of a \({\rm Binomial}(N,1/2)\) variable is at most \(C_0N^{-1/2}\).  To see this without invoking an external result, for \(N=2q\) put \(a_q={2q\choose q}4^{-q}\).  Since
\[
 \frac{a_{q+1}}{a_q}=\frac{2q+1}{2q+2},
 \qquad
 \left(\frac{2q+1}{2q+2}\right)^2\leq\frac{q+1}{q+2},
\tag{20}
\]
induction gives \(a_q\leq(q+1)^{-1/2}\), and the adjacent even case bounds the odd case.  The interval \(|M/N-1/2|\leq h\) contains at most \(2Nh+2\) integer values, so
\[
 \Pr\{|M/N-1/2|\leq h\mid N\}
 \leq C_1(\sqrt N\,h+N^{-1/2}).
\tag{21}
\]
Jensen's inequality gives \(\mathbb E\sqrt{N_i}\leq\sqrt d\).  Exponential Markov applied to \(e^{-(\log2)N_i}\) shows that
\[
 \Pr(N_i<\mu_n/2)\leq e^{-c_0\mu_n}
\tag{22}
\]
for a universal \(c_0>0\).  Splitting according to the event in (22) gives
\[
 \mathbb E\{N_i^{-1/2}\mathbf1(N_i>0)\}
 \leq\sqrt{2/\mu_n}+e^{-c_0\mu_n}
 \leq C_2d^{-1/2}.
\tag{23}
\]
Combining (21)--(23) and using \(h\geq d^{-1}\) gives, uniformly over \(h\in\mathcal H_{n,d}\),
\[
 \Pr(|V_i|\leq h,N_i>0)\leq C_3\sqrt d\,h.
\tag{24}
\]

Finally choose a fixed \(K>e\) so large that
\[
 a_K:=K\log(K/e)>2/C_{\log},
\]
and put \(D_n=\lceil Kd\rceil\).  Exponential Markov's inequality gives, because \(\mu_n\leq d\),
\[
 \Pr(N_i>D_n)
 \leq\left(\frac{e\mu_n}{D_n}\right)^{D_n}
 \leq e^{-a_Kd}.
\tag{25}
\]
The union bound and \({\rm ass:log\text{-}degree}\) therefore imply
\[
 \Pr\left(\max_{1\leq i\leq n}N_i>D_n\right)
 \leq ne^{-a_Kd}leq n^{1-a_KC_{\log}}=o(1).
\tag{26}
\]

\textit{Step 4: the graph-floor pair.}
Let
\[
 g_G(x)=c_GB\sqrt{d/n}\,(x-1/2).
\tag{27}
\]
Because \(d/n\to0\), condition (8) holds for all sufficiently large \(n\) whenever \(0<c_G\leq1\).  Equations (12), (18), (19), and (27) give
\[
 \begin{aligned}
 \mathbb E\sum_{i=1}^ng_G(X_i)^2
 &=c_G^2B^2\frac dn
 \left\{\mathbb E\sum_iV_i^2\mathbf1(N_i>0)
       +\frac n4\Pr(N_i=0)\right\}\\
 &\leq c_G^2B^2\left\{\frac12+\frac d4e^{-d/2}\right\}
 \leq C_4c_G^2B^2.
\end{aligned}
\tag{28}
\]
Fix
\[
 \eta=(1-2\alpha)/2>0.
\tag{29}
\]
Since \(s_0=B/4\), (13), (14), and (28) allow \(c_G>0\) to be chosen independently of \(n,d\) so that
\[
 \operatorname{TV}(Q_G^+,Q_G^-)\leq\eta.
\tag{30}
\]
By (9) and (27), the target separation is
\[
 \Delta_G=2c_GB\sqrt{d/n}.
\tag{31}
\]

\textit{Step 5: the interior smooth pair.}
Suppose \(L>0\).  Define the explicit odd function
\[
 \phi(t)=t(1-4t^2)^4\mathbf1\{|t|<1/2\}.
\tag{32}
\]
The fourth-order zeros at the two endpoints make the extension in (32) a \(C^3(\mathbb R)\) function, and \(\phi'(0)=1\).  Write \(C_j=\|\phi^{(j)}\|_\infty\) for \(0\leq j\leq3\), and for \(h\in\mathcal H_{n,d}\) set
\[
 g_h(x)=c_ILh^3\phi\!\left(\frac{x-1/2}{h}\right).
\tag{33}
\]
Then
\[
 \|g_h^{(j)}\|_\infty=c_ILh^{3-j}C_j,
 \qquad 0\leq j\leq3.
\tag{34}
\]
Choose \(c_I C_3\leq1\).  Since \(h\leq d^{-1/2}\to0\) and \(B,L\) are fixed, (34) verifies (8) for all sufficiently large \(n\).  The bump is zero at the isolate exposure zero, and its nonzero support is contained in \(|V_i|\leq h\).  Thus (13), (24), and (33) imply
\[
 H^2(Q_h^+,Q_h^-)
 \leq C_5c_I^2n\sqrt d\,L^2h^7.
\tag{35}
\]
Put
\[
 h_0=(L^{-2}/(n\sqrt d))^{1/7}.
\tag{36}
\]
When \(d^{-1}<h_0<d^{-1/2}\), substitution of \(h_0\) in (35), followed by a further fixed decrease of \(c_I\), yields
\[
 \operatorname{TV}(Q_{h_0}^+,Q_{h_0}^-)\leq\eta.
\tag{37}
\]
Equations (9), (33), and (36) give the separation
\[
 \Delta_I=2c_ILh_0^2
 =2c_IL^{3/7}(n^2d)^{-1/7}.
\tag{38}
\]

\textit{Step 6: the degree-lattice pair.}
Still suppose \(L>0\), put \(w_n=(4D_n)^{-1}\), and define
\[
 g_D(x)=c_DLw_n^3\phi\!\left(\frac{x-1/2}{w_n}\right),
\tag{39}
\]
where \(c_DC_3\leq1\).  The analogue of (34), together with \(D_n\to\infty\), verifies (8) for large \(n\).  If \(1\leq N_i\leq D_n\), every noncentral point on the exposure lattice obeys
\[
 |V_i|\geq
 \begin{cases}N_i^{-1},&N_i\text{ is even},\\(2N_i)^{-1},&N_i\text{ is odd},\end{cases}
 \geq(2D_n)^{-1}>w_n,
\tag{40}
\]
while oddness gives \(g_D(1/2)=0\).  Also \(g_D(0)=0\) for isolates.  Hence on
\[
 E_n=\{\max_iN_i\leq D_n\}
\]
the complete observed outcomes under the two signs in (7) both equal \(\xi_i\).  Their restrictions to \(E_n\) are identical, so (26) gives
\[
 \operatorname{TV}(Q_D^+,Q_D^-)
 \leq\Pr(E_n^c)=o(1).
\tag{41}
\]
In particular, this total variation is eventually at most \(\eta\).  On the other hand, (9) and (39) give
\[
 \Delta_D=2c_DLw_n^2\asymp Ld^{-2},
\tag{42}
\]
where both comparison constants depend only on the fixed choice of \(K\).

\textit{Step 7: reductions to risk and honest expected length.}
For any one of the preceding pairs, abbreviate the target values by \(\theta^+,\theta^-\), the data laws by \(Q^+,Q^-\), and their separation by \(\Delta\).  Given an arbitrary estimator \(T\), let
\[
 A_T=\{|T-\theta^+|\geq\Delta/2\}.
\]
On \(A_T^c\), the triangle inequality implies \(|T-\theta^-|\geq\Delta/2\).  Therefore
\[
 \begin{aligned}
 \mathbb E_{Q^+}|T-\theta^+|+\mathbb E_{Q^-}|T-\theta^-|
 &\geq\frac\Delta2\{Q^+(A_T)+Q^-(A_T^c)\}\\
 &\geq\frac\Delta2\{1-\operatorname{TV}(Q^+,Q^-)\}.
\end{aligned}
\tag{43}
\]
Thus the larger risk is at least \(\Delta(1-\eta)/4\) whenever the pairwise total variation is at most \(\eta\).

Now let \(\mathcal C\) be any interval with coverage at least \(1-\alpha\) under every law in \(\mathcal M_{n,d,L}\).  Total variation and the union bound give
\[
 \begin{aligned}
 Q^+\{\theta^+,\theta^-\in\mathcal C\}
 &\geq Q^+\{\theta^+\in\mathcal C\}
      +Q^+\{\theta^-\in\mathcal C\}-1\\
 &\geq(1-\alpha)+(1-\alpha-\operatorname{TV}(Q^+,Q^-))-1\\
 &\geq 1-2\alpha-\eta=(1-2\alpha)/2.
\end{aligned}
\tag{44}
\]
On this event the interval has length at least \(\Delta\), whence
\[
 \mathbb E_{Q^+}\operatorname{length}(\mathcal C)
 \geq\Delta(1-2\alpha)/2.
\tag{45}
\]
Taking the infima in \({\rm def:minimax\text{-}criteria}\), equations (30)--(31), (37)--(38), and (41)--(42) show that both minimax criteria are bounded below, up to fixed positive constants, by the graph term \(\sqrt{d/n}\), by the interior term \(L^{3/7}(n^2d)^{-1/7}\) when (36) is interior, and by the lattice term \(Ld^{-2}\), respectively.

\textit{Step 8: comparison with the exact supported envelope.}
Let
\[
 a_n=(n\sqrt d)^{-1/2}.
\]
For \(L>0\), (36) implies
\[
 Lh_0^2=a_nh_0^{-3/2}=L^{3/7}(n^2d)^{-1/7}.
\tag{46}
\]
If \(h_0\leq d^{-1}\), then \(L^2\geq d^{13/2}/n\), and hence
\[
 a_nd^{3/2}=d^{5/4}n^{-1/2}\leq Ld^{-2}.
\]
Evaluating the infimum in \({\rm def:frontier\text{-}envelope}\) at \(h=d^{-1}\) gives
\[
 \inf_{h\in\mathcal H_{n,d}}\{Lh^2+a_nh^{-3/2}\}
 \leq2Ld^{-2}.
\tag{47}
\]
If \(d^{-1}<h_0<d^{-1/2}\), evaluating at \(h_0\) and using (46) gives
\[
 \inf_{h\in\mathcal H_{n,d}}\{Lh^2+a_nh^{-3/2}\}
 \leq2L^{3/7}(n^2d)^{-1/7}.
\tag{48}
\]
If \(h_0\geq d^{-1/2}\), then \(L^2\leq d^3/n\), so \(Ld^{-1}\leq\sqrt{d/n}\); evaluation at \(h=d^{-1/2}\) gives
\[
 \inf_{h\in\mathcal H_{n,d}}\{Lh^2+a_nh^{-3/2}\}
 \leq2\sqrt{d/n}.
\tag{49}
\]
When \(L=0\), the infimum is attained at \(d^{-1/2}\) and equals \(\sqrt{d/n}\).

In the regime of (47), combine the graph and lattice pairs; in the regime of (48), combine the graph and interior pairs; and in the regime of (49), or when \(L=0\), use the graph pair.  Because a minimax criterion bounded below by \(c_1x\) and by \(c_2y\) is bounded below by \(\min(c_1,c_2)(x+y)/2\), equations (31), (38), (42), and (47)--(49) yield, with a fixed \(c>0\),
\[
 R^*_{n,d,L}\geq c\left[\sqrt{d/n}+
 \inf_{h\in\mathcal H_{n,d}}\{Lh^2+(n\sqrt d\,h^3)^{-1/2}\}\right],
\tag{50}
\]
and the same inequality with \(R^*_{n,d,L}\) replaced by \(\Lambda^*_{n,d,L}(\alpha)\).  The bracket in (50) is exactly \(\mathfrak r^{\mathrm{sup}}_{n,d,L}\) by \({\rm def:frontier\text{-}envelope}\).  All constants used above depend only on the fixed class constants and \(\alpha\), never on \(n\) or \(d\).  This proves the asserted unrestricted converse.